\documentclass[11pt]{article}
\usepackage[a4paper,margin=2.5cm]{geometry}
\usepackage{amsmath,amssymb,booktabs,xcolor,enumitem}
\newcommand{\ok}{\textcolor{green!40!black}{\textbf{VERIFIED}}}
\newcommand{\bad}{\textcolor{red!70!black}{\textbf{COUNTEREXAMPLE}}}
\newcommand{\gap}{\textcolor{orange!80!black}{\textbf{PROOF GAP}}}
\newcommand{\fix}{\textcolor{blue!70!black}{\textbf{LOCAL REPAIR}}}
\newcommand{\todo}{\textcolor{gray!70!black}{\textbf{PENDING}}}
\newcommand{\lean}[1]{\texttt{#1}}
\title{Validation ledger for Todorcevic--Tyros, A Dual Ramsey Theorem for Trees}
\author{lean-dualtree}
\date{\today}

\begin{document}
\maketitle

\section{Policy}
This ledger tracks arXiv:2207.14599v1 statement by statement.  We distinguish
false printed statements, type/domain errors, local repairs, missing arguments,
and statements actually validated in Lean.  A corrected statement is never
silently substituted for the printed one: the repair is recorded here first.

\section{Audit summary}
\begin{center}
\begin{tabular}{p{.22\linewidth}p{.18\linewidth}p{.52\linewidth}}
\toprule
Location & Status & Finding \\ \midrule
\S3.1 auxiliary order & \bad &
Reverse lexicographic tie-breaking is incompatible with the claimed canonical
order-preserving isomorphism.  A three-node binary witness is formalized. \\
Remark 2 & \bad &
With the printed order, the interior of a semi-complete skew tree need not be
skew.  Finite tests support forward lexicographic tie-breaking. \\
Remark 1 & \bad &
There is an off-by-one domain error, and the later extension assertion has a
separate finite counterexample even after the forward-lex repair. \\
Theorem 3, $\ell=1$ & \bad &
Evaluated spans collapse for a one-letter alphabet, so span inclusion no
longer expresses syntactic refinement. \\
Corollary 22 & \bad &
The displayed numerical threshold fails in the smallest starred example. \\
Theorem 25 comparison & \bad &
The displayed reverse inequality from TGR to PTGR fails in a binary example. \\
Remark 3(ii) & \bad/\fix &
Union-insensitivity fails for disjoint sets, but is valid for overlapping
sets; the application uses overlapping consecutive pairs. \\
Lemma 16 / Claim 17 & \gap &
The recursive block budget and dimension indices do not match Corollary 9. \\
Theorem 13 gluing & \fix/\gap &
Several complements and tail domains are mistyped; the induction architecture
still looks plausible. \\
Lemma 27 & \gap &
Main bottleneck: projection domains, auxiliary-variable coding, extension of
the coloring, and preservation of variable fibres all need repair. \\
Lemma 28 & \gap &
Depends on Lemma 27 and suppresses signature transport during iteration. \\
Theorem 25 induction & \todo &
The pullback coloring must be proved simple before the induction hypothesis is
applied.  Finite tests are positive. \\
Theorem 29 & \todo &
The product theorem is stated without a separate proof. \\
\bottomrule
\end{tabular}
\end{center}

\section{Detailed issues}

\subsection{The auxiliary node order}
The paper orders shorter nodes first and uses reverse lexicographic order on
ties.  For
\[
S=\{\varnothing,0,10\},
\]
the canonical tree/lexicographic map from $2^{<2}$ sends
$0\mapsto0$ and $1\mapsto10$.  In the source the printed tie-break has
$1\preccurlyeq0$, while in the target the length clause gives
$0\prec10$.  Thus the canonical map cannot preserve the printed order.

\paragraph{Lean.} \ok
\lean{DualTree.OrderAudit.canonical_order_counterexample}.
The local forward-lex repair is checked by
\lean{DualTree.OrderAudit.forward_order_repairs_counterexample}.

\paragraph{Proposed repair.}
Use forward lexicographic order among nodes of equal length.  Exhaustive checks
found zero failures for the canonical-order and interior tests in the tested
binary and ternary ranges.  This does not by itself repair Remark 1.

\subsection{Remark 2}
With the printed order,
\[
S=\{\varnothing,0,1,10,11\}
\]
gives a semi-complete example whose interior
$\{\varnothing,1\}$ violates the one-successor direction convention.
The statement should be reproved after the order convention is corrected.

\paragraph{Lean status.}
The concrete witness is certified by the executable Section 3 checker:
\lean{DualTree.BinarySkewAudit.remark2_witness_is_semicomplete_printed} and
\lean{DualTree.BinarySkewAudit.remark2_interior_not_skew_printed}.  Equivalence
of this executable checker with the final propositional skew-tree definition
is still \todo.

\subsection{Remark 1}
Two independent problems were found.
\begin{enumerate}[label=(\alph*)]
\item The condition $m\le n-|s|$ permits concatenations reaching level $n$,
outside $b^{<n}$.  The evident safe bound is $m\le n-|s|-1$, modulo the
authors' intended indexing convention.
\item The claimed extension of every skew $T'\subseteq T$ to a complete skew
subtree of the spliced tree has a finite counterexample even with forward lex.
This needs a new statement, not just an index correction.
\end{enumerate}

\paragraph{Lean status.}
\ok The substantive forward-lex failure is now certified in
\lean{DualTree.Remark1Audit.no_forward_complete_extension}.  The chosen
spliced tree and the prescribed intersection tree are separately certified
skew by \lean{DualTree.Remark1Audit.forwardS_is_skew} and
\lean{DualTree.Remark1Audit.forwardTprime_is_skew}.  The remaining task for
this item is not verification of the printed statement---it is false---but
identifying the intended replacement needed downstream.

\subsection{Span inclusion and Theorem 3 for one-letter alphabets}
The paper defines refinement using
\[
[g]_{\Lambda}\subseteq[f]_{\Lambda}.
\]
If $|\Lambda|=1$, every variable word evaluates to the same constant word.
For $b=2$, $n=2$, $k=m=1$, coloring by whether the unique variable root is
$\varnothing$ gives two colors, while every one-variable word is a refinement
of every candidate in the printed sense.  Hence Theorem 3 is false literally
for $\ell=1$.

\paragraph{Repair and Lean status.}
\ok The algebraic issue is now completely separated from the tree geometry.
The singleton-alphabet collapse is formalized as
\lean{DualTree.SpanAudit.singleton_span_subset}.  Conversely, if the alphabet
contains two distinct letters and every source variable occurs, then
\lean{DualTree.SpanAudit.span_subset_iff_substitution} proves that evaluated
span inclusion is equivalent to uniform syntactic substitution.  The proof is
factored through preservation of source constants and uniformity on each
source variable fibre.

Thus later uses of span inclusion are justified once the paper is restricted
to $|\Lambda|\ge2$ (or, preferably, once refinement is defined syntactically).
The remaining tree-specific task is to combine this algebraic lemma with the
variable-support/root conditions.

\subsection{Corollary 22}
Take
\[
d_0=d_1=0,\quad d_2=1,\quad b=\ell=r=2,\quad k=n=1,
\]
and color $(w,\varnothing)$ by $w(\varnothing)$.  The marked point is forced
to be the root, so the coloring is smooth with $\Gamma_1=D_2$.  The displayed
threshold permits height one, but every genuine one-dimensional candidate has
both colors in its starred span.  Thus the numerical statement is false as
printed.

\paragraph{Lean status.}
The height-one mechanism is now source-facing.  The uniqueness of the marked
node is \lean{DualTree.Corollary22Audit.node_eq_root};
\lean{DualTree.Corollary22Audit.color_is_smooth} proves the specialized
Definition~21 smoothness condition; and
\lean{DualTree.Corollary22Audit.no_monochromatic_complete_candidate} proves
that no complete one-variable candidate has monochromatic starred span.
Moreover,
\lean{DualTree.Corollary22Audit.canonical_good_for_every_coloring}
proves that at ambient height one the canonical one-dimensional subspace is
$c$-good for every coloring, i.e. the local fact yielding the relevant
MTHJ value $1$.  The equally trivial numerical identity
$\operatorname{CT}(1,1,2,1,r)=1$ has not yet been wrapped in a general CT
definition in Lean.

There is also a boundary-definition issue: the corollary permits $d_2=0$ but
uses $\operatorname{CT}(1,k,b,d_2,r)$, while Theorem 10 defines CT only for
positive vector dimension.

Finite search at the next binary height found no analogous obstruction, so a
small dimensional margin may repair the argument.

\subsection{The comparison after Theorem 25}
The displayed estimate
\[
\operatorname{PTGR}(l,m,b,r)
\le \operatorname{TGR}(k_l+1,m,b,\ell,r)
\]
fails already for $b=\ell=r=2$, $l=1$, $m=2$.  The right side is $2$, while a
simple coloring at ambient height two defeats the sole height-two candidate.
The qualitative reduction may survive with different parameters.

The notation $\operatorname{PTGR}(l,m,b,r)$ also suppresses the fixed alphabet
size $\ell$, although Lemmas 27 and 28 carry $\ell$ explicitly.  We should use
$\operatorname{PTGR}(l,m,b,\ell,r)$ in the formal development.

\subsection{Definition 24}
$I_S$ was defined for complete skew trees, but Definition 24 applies it to a
semi-complete tree.  The intended canonical transport should be defined
explicitly rather than overloaded.

\subsection{Lemma 16 and Claim 17}
Corollary 9 returns a block length.  The recursion
\[
q_p=Q(\ell,d_2^p,d_2q_{p+1},r)
\]
is then used as though it already contains room for $q_{p+1}$ blocks.
A safe recursion is of the form
\[
q_p=q_{p+1}Q(\ell,d_2^p,d_2q_{p+1},r).
\]
Claim 17 also shifts the $M_p$ dimensions by one.  We will first formalize the
qualitative simultaneous refinement and only then restore quantitative bounds.

\subsection{Theorem 13 gluing}
The double induction itself appears well founded, but the final gluing contains
concrete source defects:
\begin{itemize}
\item the definition of $B$ uses the wrong primed set in its cardinality
condition;
\item $\overline Y_i$ is written as a complement in $[b^k]$ although
$Y_i\subseteq[b^M]$;
\item the tail formulas quantify over $b^{n-M}$ where the word domain is
$b^{<n-M}$.
\end{itemize}
These are candidates for local source repairs, but the complete gluing should
be Lean-checked before being marked harmless.

\subsection{Lemma 27}
This is the present structural bottleneck.
\begin{enumerate}[label=(\alph*)]
\item With $n'=n-m'$, the map
$P_i(z)=I_T(I_T^{-1}(t_i)^\frown z)$ can leave the domain of $I_T$ for
frontier nodes at the larger intrinsic height.  A common safe tail loses one
level, or the cones need different tail heights.
\item The requested surjection from $\{0,\ldots,m'\}$ to all interior
variables can be cardinally impossible and can assign incomparable variables
to a cone.  Only ancestor variables visible from the given frontier cone can
be legitimate targets.
\item The enlarged alphabet
$\Lambda'=\Lambda\sqcup\{0,\ldots,m'\}$ is introduced but is not consistently
used in the mixed-product domain.
\item The induced coloring is defined on the starred part, while Corollary 22
takes a coloring of the full mixed-product space.  A compatible extension is
not supplied.
\item The reconstruction tests the occurrence position $P_i(z)$ rather than
the root $P_i(y)$ of the variable $v_y$, so different occurrences of one
variable fibre can be frozen differently.
\end{enumerate}
Therefore the final inclusion of every continuation of the fixed signature in
the monochromatic image is not established.  We should reconstruct the local
coding lemma rather than formalize these lines verbatim.

\subsection{Lemma 28 and Theorem 25}
Lemma 28 iterates Lemma 27 and additionally needs explicit transport between a
model signature in $b^{<m}$ and the corresponding signature observable in the
current variable word.

In Theorem 25, after pulling the signature-coloring back along $Q$, the proof
invokes the induction hypothesis without showing that the new coloring is
simple.  Tree-only finite tests are positive: the modeled $Q$ is onto in the
tested $l=1,2$ cases and tested simple colorings remain simple.  This should be
isolated as a lemma.

\subsection{Appendix Remark 3(ii)}
The unrestricted union assertion is false for disjoint $L_1,L_2$.  It becomes
true when $L_1\cap L_2\ne\varnothing$, which is sufficient for the actual
application $L_j=\{a_j,a_{j+1}\}$.

\paragraph{Lean status.} \ok
Definition~30's exact assignment relation is formalized as
\lean{DualTree.StarInsensitivity.StarRelated}, and
\lean{DualTree.StarInsensitivity.starRelated_iff} identifies it with the
three printed clauses.  The unrestricted printed claim is refuted by the
kernel-checked
\lean{DualTree.StarInsensitivity.disjoint_union_failure}.
The corrected overlap statement actually used in the proof is
\lean{DualTree.StarInsensitivity.starInsensitive_union_of_overlap}.
Its proof normalizes both assignments through a common bridge letter in
$L_1\cap L_2$, making the suppressed transitivity explicit.

\subsection{Lemma 31}
Checking the rendered PDF removes two OCR false alarms: the lemma assumes
$k\ge2$, and the proof later says to assume $k\ge3$.  Two genuine index shifts
remain:
\begin{itemize}
\item at stage $i$ the new variable is in slot $i-1$, not $i$;
\item the displayed cardinality of $J$ has $m-i-1$, while substitution into
$d_{i-1}-p_i$ gives the term with $m-i$.
\end{itemize}
These currently look local.

\subsection{Theorem 29}
The product theorem is only asserted to follow by modifying the preceding
argument.  It remains unverified.

\section{Computational evidence retained}
Before Lean formalization, exhaustive finite programs established or tested:
\begin{itemize}
\item failures of the printed auxiliary order and Remark 2, with zero failures
of the corresponding local tests after the forward-lex change;
\item a separate forward-lex counterexample to Remark 1's extension claim;
\item collapse of refinement for $|\Lambda|=1$ and no nonsyntactic inclusion
in the tested two-letter instances;
\item the height-one Corollary 22 counterexample;
\item the small counterexample to the displayed PTGR/TGR inequality;
\item 1244 finite overlap tests supporting the repaired Remark 3(ii);
\item onto behavior of the Theorem 25 signature map and simplicity preservation
in the tested tree-only cases.
\end{itemize}

\section{Formalization order}
\begin{enumerate}
\item finite homogeneous trees, prefix and the two candidate auxiliary orders;
\item skew, complete skew and semi-complete trees, followed by Remarks 1--2;
\item variable words, evaluation, syntactic substitution and span inclusion;
\item the $|\Lambda|\ge2$ substitution/refinement lemma and the unary
counterexample;
\item starred insensitivity and the Appendix argument;
\item mixed products and Theorem 13;
\item signatures and a repaired Lemma 27;
\item Lemma 28 and Theorem 25;
\item Theorem 3 and finally the product theorem.
\end{enumerate}

Qualitative statements take priority over the large recursive bounds until the
underlying constructions are type-correct.

\end{document}
